% =====================================================================
%  Chapter 12 — Principle: Connection (Speak Across Every Boundary)
%  Map: research/chapter-map.md (Chapter 12)
%  Shape: idea → origin → cost → what it opens up. Ideas, not how-to.
%  Rule: the author is never quoted in body text; short quotes only in
%  footnotes. Full statements live in interviews/.
% =====================================================================
\chaptersubtitle{Speak Across Every Boundary}
\chapter{Connection}
\label{ch:transport}

Moving data between devices is as basic to modern computing as arithmetic.
Yet most languages leave it to libraries, each with its own way of packaging
data, its own ways of failing and its own security holes. Silica treats the
movement of data between machines as a full and equal part of the language
itself, understood by the compiler and checked like everything
else.\footnote{Specification §1.2,
\src{silica/compiler/design_documents/silica-specification.md}, line 28.}

\section{The long road to the network}

This principle has the longest history of any in the book. SwErl's unfinished
work of 2024 could already let a Swift program join a cluster of Erlang
machines. BEAM\_SEMP's work of 2025 made such connections trustworthy: every
peer had to prove its identity, only known peers were admitted, and each
connection was handled on its own. Both ideas reappear, almost unchanged, in
Silica's designs.

\section{In, out, and in again}

Inside Silica, the road was not straight. In September 2026 networking was
made part of the language. Within days that was reversed. Then, in November, a
third position emerged, and it is the principle. The plumbing of networks---the
operating system's connections and ports---is not the language's to own. But
\emph{what} travels between machines, how it is packaged, and how it is checked
when it arrives are questions about data and meaning, which is to say questions
for the language.

% TODO-INTERVIEW: The three changes of position: why not networking, why not
%   a library, and why is moving data different from plumbing?

\section{Trust nothing that arrives}

Silica treats every byte that arrives from elsewhere as untrusted. Checking and
unpacking are the same act: data is either fully valid or it is refused, never
half-accepted. The ways data can be packaged and carried form a fixed set
defined by the language, so no program can quietly invent a weaker one. Adding
a new way is a change to the language, examined like any other.

The same rule covers data a program stored itself. Silica's local persistent
storage uses the same transport, so a file read back from a disk is checked
exactly as a message from another machine would be: something may have changed
it since it was written (Chapter~\ref{ch:memory}).

The simplest connection Silica offers is documented as exactly that: no
encryption, no identity. Secure exchange is being built on top of it, openly.

\section{Speaking an old friend's language}

The first way Silica packages data is Erlang's own. Silica machines can
exchange data with Erlang systems, and Project BEES---an effort to let programs
written in Erlang, Elixir, Gleam and LFE be compiled to Silica---had two Silica
machines speaking Erlang's distribution protocol by November 2026. BEES's rule
is modest and telling: make it possible, not easy, and ask Silica to change
only where something would otherwise be impossible. The trust layer first built
for BEAM\_SEMP is being rebuilt there, on Silica.

% TODO-INTERVIEW: How much did BEES shape Silica's approach to connection?
%   The dates suggest it did; no document says so.

\section{What it costs}

Making connection part of the language means every new way of packaging or
carrying data must earn its place in the language itself, which is slower than
writing a library. And refusing half-valid data means some messages other
systems would tolerate are turned away.

\section{What it opens up}

\begin{itemize}
  \item \textbf{One design across many machines.} The same actors, with the
        same kinds of data, can run on a sensor, a gateway and a server, with
        only the connection between them changing.
  \item \textbf{A safer edge.} Arriving data is always checked by routines
        the compiler generates from the type of data the program expects,
        not by hand-written parsing code, so a whole family of attacks on
        data parsing has nowhere to start.
  \item \textbf{A bridge to a living ecosystem.} Through Erlang's own format,
        decades of Erlang and Elixir systems can talk to Silica, and through
        BEES their programs may one day run on machines the Erlang virtual
        machine cannot reach.
  \item \textbf{Fleets that are trustworthy by default.} Devices in the field
        can be built to admit only known peers, the lesson of BEAM\_SEMP made
        part of the foundation.
\end{itemize}

\section{Where the principles meet}

Connection is where the book's principles come together: independent actors
(Chapter~\ref{ch:actors}), effects stated where they happen
(Chapter~\ref{ch:effects}), danger made visible (Chapter~\ref{ch:danger}) and one
meaning on every machine (Chapter~\ref{ch:machines}). A program that keeps all of
those promises on one machine can keep them across many.

\section{Beyond the operating system}

Asked what all of this is finally for, Barney points past programming
languages altogether, to the way our devices are organised. They are, in
Barney's view, a mess. The same word processor is installed again and again, on
every machine a person owns. Every one of those machines carries its own
operating system, an arrangement whose basic design, as Barney sees it, was set
in the late 1950s and early 1960s, and each wastes space and attention on its own
copy. Barney's verdict on decades of competition between operating systems is
that nobody won: not the users, and not the companies that made
them.\footnote{``The operating system wars are over. Everybody lost.''
Interview, 2026-11-13; also the opening line of the Organism draft.} It is time
to move on.

What moving on might look like is sketched, in very early form, in a draft
specification Barney calls Organism.\footnote{\src{https://github.com/yenrab/organism},
Draft Specification v0.3.} In it, a person, or a family or a team, has one
persistent computational life, and every device they use is simply an organ of
that life. No single machine is in charge. Actors take the place of processes
and threads, and every interaction is a message. Authority comes only from
capabilities, never from ambient permissions. Keeping things alive across
restarts, moves between devices and even the loss of a device is the system's
job, not each application's. Files give way to a shared, content-addressed
memory. A user interface is a temporary projection of actors that can appear on
whichever device is at hand. Failure is treated as normal, something to
contain and heal rather than to fear.

Read against this book, Organism is the ten principles at the scale of a
person's whole digital life. Independent actors are its basic unit
(Chapter~\ref{ch:actors}). Authority you can see is its security model
(Chapter~\ref{ch:danger}). Failing gently is its everyday behaviour
(Chapter~\ref{ch:failsoft}). And connection across devices that sleep, wake and
roam is its reason to exist.

Organism is not built. It is an imagining, and its own pages list the research
it still needs. But Barney sees no reason something like it could not be built
in Silica, and hopes to begin soon. If that happens, a language that started on
24~December 2025 with the realisation that its creator was solving the wrong
problem will have found the right one.

% TODO-INTERVIEW: Organism and SEMP — how do SEMP's churn ideas feed Organism?
